%=============================================================
% CHAPTER 2 — SOFTWARE REQUIREMENT SPECIFICATIONS (SRS)
%=============================================================

\chapter{Software Requirement Specifications (SRS)}

Getting requirements right before writing a single line of code is one of those lessons
that takes most developers a bad experience to learn. This chapter documents what
MediLingo AI is actually supposed to do: the problem motivating it, the boundaries it works
within, and the complete set of functional and non-functional requirements. Use cases and
the domain model sit at the end, building on everything before them. The system design
in Chapter~3 should be fully traceable back to this chapter.

%=============================================================
\section{Introduction}
%=============================================================

%-------------------------------------------------------------
\subsection{Problem Statement}
%-------------------------------------------------------------

Pakistan has one of the largest Urdu-speaking populations in the world, yet there is no
publicly accessible, AI-driven medical tool that accepts queries in Urdu or Roman Urdu and
grounds its answers in verified medical documents. The gap has two distinct layers.

The first is language. Every AI medical tool of any quality is built for English. They
cannot process a query written in Urdu, and even when a user translates their question
manually, there is no mechanism to return a response in their native language. The second
layer is reliability. General-purpose LLMs that do support Urdu produce answers from
training weights, not from a curated clinical corpus. In healthcare, a confident but wrong
answer is not just unhelpful --- it carries real risk. Medical facts, dosage details and
symptom descriptions need to come from somewhere traceable, not from statistical patterns
baked into model parameters.

MediLingo AI addresses both. It accepts queries in English, Urdu and Roman Urdu, and every
response is grounded in passages retrieved from an indexed medical knowledge base at query
time --- not recalled from model memory.

%-------------------------------------------------------------
\subsection{Motivation}
%-------------------------------------------------------------

The case for building this came down to three things that kept coming back during the
requirements phase.

\begin{itemize}
  \item \textbf{Language accessibility:} Forcing an Urdu speaker to translate their medical
        question before they can get a useful answer is not a minor inconvenience. For a lot
        of people it is a genuine barrier. Removing the translation step lowers the cost of
        accessing health information for tens of millions of Urdu speakers. That felt like
        a concrete, achievable goal rather than an abstract one.

  \item \textbf{Factual reliability:} The hallucination problem is well-known. RAG is not
        a complete solution, but it meaningfully narrows the gap by tying every response
        to specific retrieved passages from a known knowledge base. At minimum, it makes
        the system's outputs traceable to a source document --- something a standard LLM
        pulling from training memory alone cannot do.

  \item \textbf{Practical feasibility:} The components needed for this --- open-source
        embedding models, a cloud-hosted LLM, a local vector database --- are mature enough
        now that a working system is achievable without proprietary data or specialised
        hardware. That was not obviously true two or three years ago.
\end{itemize}

%-------------------------------------------------------------
\subsection{Objectives}
%-------------------------------------------------------------

The primary objectives of the MediLingo AI project are:

\begin{itemize}
  \item Build a web-based medical chatbot accepting queries in English, Urdu and Roman Urdu,
        returning responses in the same language as the query
  \item Ground every response in verified medical PDF documents through a RAG pipeline ---
        the LLM synthesises and explains, but does not recall clinical facts from memory
  \item Give users a persistent, session-based chat interface so that separate medical
        conversations can be maintained and revisited over time
  \item Implement JWT dual-token authentication with HTTP-only refresh cookies and silent
        renewal so the user's session does not interrupt them
  \item Integrate voice input and output through the Web Speech API, making the system
        accessible to users who prefer speech interaction
  \item Provide an admin interface for user account management and knowledge base
        re-ingestion from PDF files placed in the server datasets directory
\end{itemize}

%-------------------------------------------------------------
\subsection{Scope}
%-------------------------------------------------------------

\textbf{In scope:}
\begin{itemize}
  \item Web application accessible from any modern browser
  \item Query input and response output in English, Urdu and Roman Urdu
  \item RAG pipeline using locally ingested medical PDF documents as the knowledge base
  \item User registration, login, logout and password reset
  \item Chat session creation, continuation, deletion and history viewing
  \item Voice input (speech-to-text) and voice output (text-to-speech)
  \item Admin panel for user management and PDF ingestion trigger
  \item Containerised deployment via Docker and Docker Compose
\end{itemize}

\textbf{Out of scope:}
\begin{itemize}
  \item Native mobile application (iOS or Android)
  \item Diagnostic conclusions or prescriptive medical advice --- the system provides
        information only, not clinical diagnosis
  \item Real-time consultation with a licensed medical professional
  \item Support for languages other than English, Urdu and Roman Urdu
  \item Integration with external electronic health record (EHR) systems
\end{itemize}

%-------------------------------------------------------------
\subsection{Constraints}
%-------------------------------------------------------------

\begin{itemize}
  \item A valid \texttt{GROQ\_API\_KEY} environment variable is required; the backend
        starts successfully without it but every query returns an error until it is set
  \item Voice input and output depend on the Web Speech API, which has full support only
        in Chromium-based browsers (Chrome, Edge) and partial support in Firefox; Safari
        is unsupported for voice input
  \item ChromaDB is used in embedded mode, meaning it is co-located with the backend
        process and does not support distributed or replicated vector storage out of the box
  \item The \texttt{sentence-transformers} model (\texttt{all-MiniLM-L6-v2}) is loaded
        fully into RAM on startup; the server requires a minimum of 512 MB of available
        memory for the ML components alone
  \item Medical knowledge is limited to the content of PDF files placed in
        \texttt{backend/datasets/} and ingested through the admin pipeline; the system
        cannot answer questions on topics not covered by those documents
  \item The language-prefix mechanism relies on the Groq-hosted LLaMA 3.3-70b model's
        built-in multilingual capability; the quality of Urdu and Roman Urdu responses
        depends on the model's training data and may vary by medical sub-domain
\end{itemize}

%-------------------------------------------------------------
\subsection{Assumptions}
%-------------------------------------------------------------

\begin{itemize}
  \item All medical PDF documents provided for ingestion are written in English; the
        embedding model operates on English text, and translating source documents is
        outside the scope of the system
  \item Users have a stable internet connection sufficient for REST API calls and Groq
        inference requests
  \item The Groq API service remains available during system operation; no fallback LLM
        is configured
  \item Users interact with the system through a web browser on a desktop or laptop
        device; mobile browser support is functional but not the primary design target
  \item The admin user is technically familiar enough to place PDF files in the correct
        server directory and invoke the ingestion endpoint
\end{itemize}

%-------------------------------------------------------------
\subsection{System Workflow Diagram}
%-------------------------------------------------------------

Figure~\ref{fig:workflow} shows what actually happens to a query from the moment the user
hits submit to the moment the response lands on screen. The language prefix gets stripped
early in the process so that vector search always runs against clean medical text ---
this is the key design detail that makes multilingual retrieval work correctly.

\begin{figure}[H]
  \centering
  \includegraphics[width=0.55\textwidth]{workflow}
  \caption{MediLingo AI --- End-to-End System Workflow for Multilingual Medical Query}
  \label{fig:workflow}
\end{figure}

%=============================================================
\section{Overall Description}
%=============================================================

%-------------------------------------------------------------
\subsection{Product Perspective}
%-------------------------------------------------------------

MediLingo AI is a self-contained web application. It does not plug into any larger platform
and depends on only three external services: the Groq API for LLM inference, MongoDB for
persistent storage, and ChromaDB for document retrieval. All three can run on the same host
as the FastAPI backend or on separate infrastructure without any changes to the application
code. During development, a Vite proxy forwards all API traffic from the browser to
\texttt{localhost:8000}, which means browser cookies work without fighting CORS --- one
fewer thing to debug during development.

%-------------------------------------------------------------
\subsection{User Characteristics}
%-------------------------------------------------------------

\textbf{Regular User:}
Someone looking for medical information --- could be fluent in English, Urdu or Roman Urdu,
or switching between all three. The interface requires basic internet literacy and nothing
more. Users should not need to understand anything about RAG or how the backend works. Voice
support is there specifically for people who have difficulty typing or just prefer speaking.

\textbf{Admin User:}
Someone with enough technical familiarity to place PDF files in the correct server directory
and trigger the ingestion endpoint from the admin panel. They authenticate through a
dedicated admin login route that verifies the role field in the JWT payload before granting
access to admin-only endpoints.

%=============================================================
\section{Specific Requirements}
%=============================================================

%-------------------------------------------------------------
\subsection{External Interface Requirements}
%-------------------------------------------------------------

\textbf{User Interface:}
A React SPA delivered in the browser. The interface includes a landing page, auth pages
(login, register, forgot password, reset password), a chat page with a collapsible session
sidebar, a profile page and an admin panel. No plugin or browser extension is required.

\textbf{Hardware Interface:}
A standard desktop or laptop with a microphone for voice input, and speakers or
headphones for voice output. No specialised hardware needed.

\textbf{Software Interface:}
The backend exposes a documented REST API at \texttt{/docs} (Swagger UI). The frontend
uses this interface exclusively. The backend communicates with MongoDB via Motor and with
ChromaDB via its Python client.

\textbf{Communication Interface:}
All frontend-to-backend communication is HTTP/HTTPS with JSON bodies. Auth credentials go
as a Bearer token in the \texttt{Authorization} header. Refresh tokens are transmitted as
HTTP-only cookies.

%-------------------------------------------------------------
\subsection{Functional Requirements}
%-------------------------------------------------------------

\renewcommand{\arraystretch}{1.35}
\begin{longtable}{|>{\bfseries\color{navyblue}\centering}p{1.3cm}|p{4.2cm}|p{7.8cm}|}
\caption{Functional Requirements of MediLingo AI}
\label{tab:fr} \\
\hline
\theader{ID} & \theader{Requirement} & \theader{Description} \\
\hline
\endfirsthead
\hline
\theader{ID} & \theader{Requirement} & \theader{Description} \\
\hline
\endhead
\hline
\endfoot

FR-01 & User Registration &
  The system shall allow a new user to create an account using a display name, email
  address and password. Duplicate email addresses must be rejected. \\
\hline
FR-02 & User Login &
  The system shall authenticate a registered user with their email and password, issue a
  JWT access token and set an HTTP-only JWT refresh cookie. \\
\hline
FR-03 & Silent Token Refresh &
  The frontend shall automatically call \texttt{POST /auth/refresh} when the backend
  returns HTTP 401, retry the original request with the new access token, and redirect to
  login only if the refresh token has also expired. \\
\hline
FR-04 & User Logout &
  The system shall invalidate the refresh cookie and clear in-memory auth state on
  explicit logout. \\
\hline
FR-05 & Password Reset via Email OTP &
  The system shall allow a registered user to reset their password by requesting a
  one-time passcode sent to their registered email address. \\
\hline
FR-06 & Multilingual Query Submission &
  The system shall accept medical queries in English, Urdu and Roman Urdu and return
  a response in the same language as the query. \\
\hline
FR-07 & RAG-Grounded Response Generation &
  Every query response shall be generated by first retrieving the top-K semantically
  relevant passages from the ChromaDB vector store and including them in the LLM prompt. \\
\hline
FR-08 & Language-Prefix Stripping &
  The backend shall strip the language instruction prefix from the query text before
  computing the embedding, so that vector search operates on the medical content
  alone. \\
\hline
FR-09 & Chat Session Creation &
  An authenticated user shall be able to start a new chat session. The session shall be
  saved to MongoDB and appear in the session sidebar. \\
\hline
FR-10 & Chat Session Management &
  An authenticated user shall be able to view, continue and delete any of their saved
  chat sessions. All messages within a deleted session shall be removed. \\
\hline
FR-11 & Message Persistence &
  Every user query and assistant response within a session shall be saved as a
  \texttt{MessageDoc} in MongoDB so that session history survives page reload. \\
\hline
FR-12 & Voice Input &
  The system shall provide a microphone button that activates the browser's Web Speech
  API to transcribe speech into the query input field. \\
\hline
FR-13 & Voice Output &
  The system shall read assistant responses aloud using the browser's SpeechSynthesis API
  after each response is received. \\
\hline
FR-14 & User Profile Management &
  An authenticated user shall be able to view and update their display name and email
  address from the profile page. \\
\hline
FR-15 & Admin: User Management &
  An admin user shall be able to view all registered accounts and activate or deactivate
  individual users. \\
\hline
FR-16 & Admin: PDF Ingestion &
  An admin user shall be able to trigger a synchronous or background ingestion of PDF
  files from \texttt{backend/datasets/} into the ChromaDB vector store via the admin
  panel. \\
\hline
\end{longtable}

%-------------------------------------------------------------
\subsection{Non-Functional Requirements}
%-------------------------------------------------------------

\renewcommand{\arraystretch}{1.35}
\begin{longtable}{|>{\bfseries\color{navyblue}\centering}p{1.3cm}|p{3.5cm}|p{8.5cm}|}
\caption{Non-Functional Requirements of MediLingo AI}
\label{tab:nfr} \\
\hline
\theader{ID} & \theader{Category} & \theader{Requirement} \\
\hline
\endfirsthead
\hline
\theader{ID} & \theader{Category} & \theader{Requirement} \\
\hline
\endhead
\hline
\endfoot

NFR-01 & Performance &
  End-to-end query response time (from submission to display) shall not exceed 15 seconds
  under normal Groq API latency conditions. \\
\hline
NFR-02 & Security &
  User passwords shall be stored as bcrypt hashes. Access tokens shall never be persisted
  to localStorage or sessionStorage. Refresh tokens shall be stored in HTTP-only, Secure
  cookies only. \\
\hline
NFR-03 & Reliability &
  The RAG engine and ingestion pipeline shall be initialised as a singleton at backend
  startup. If initialisation fails, the backend shall still start and return a descriptive
  error on query endpoints rather than crashing entirely. \\
\hline
NFR-04 & Availability &
  The system shall be deployable as a set of Docker containers that can be started with a
  single \texttt{docker-compose up} command, minimising manual deployment steps. \\
\hline
NFR-05 & Maintainability &
  Backend logic shall be separated into distinct FastAPI routers
  (\texttt{auth}, \texttt{sessions}, \texttt{query}). Frontend business logic shall be
  isolated in custom React hooks, separate from UI components. \\
\hline
NFR-06 & Portability &
  The application shall run on any machine with Docker installed, regardless of the
  underlying operating system. \\
\hline
NFR-07 & Browser Compatibility &
  The web application shall function correctly on Chrome 120+, Edge 120+ and Firefox 121+.
  Voice features shall degrade gracefully (button hidden or disabled) on unsupported
  browsers. \\
\hline
NFR-08 & Usability &
  The language selector, voice button and session sidebar shall be visible and operable
  without reading any documentation. The interface shall be fully functional on a
  1280$\times$720 desktop viewport. \\
\hline
\end{longtable}

%=============================================================
\section{Software Product Features}
%=============================================================

MediLingo AI exposes the following top-level product features:

\begin{itemize}
  \item \textbf{Cross-Lingual Medical Chat} --- queries and responses in English, Urdu or
        Roman Urdu, with language routing handled automatically by the prefix mechanism in
        the frontend and backend

  \item \textbf{RAG-Grounded Answers} --- each response is generated from retrieved passages
        rather than from the LLM's memory; medical content is always traceable to a source
        document

  \item \textbf{Persistent Session History} --- users maintain multiple named conversations
        and can return to any of them at any time; message history survives page reloads
        because it is stored in MongoDB, not in browser memory

  \item \textbf{Voice Interaction} --- speech-to-text for query input and text-to-speech
        for response playback; both use the browser's built-in Web Speech API with no
        additional plugin required

  \item \textbf{Secure Authentication} --- 15-minute access tokens kept in memory,
        7-day refresh tokens in HTTP-only cookies, automatic silent renewal so the user's
        session does not get interrupted

  \item \textbf{Admin Knowledge Management} --- administrators can expand or refresh the
        medical knowledge base by triggering a new PDF ingestion run through the admin panel

  \item \textbf{Email-Based Password Recovery} --- users can regain access via a time-limited
        one-time passcode sent to their registered email address
\end{itemize}

%=============================================================
\section{Software System Attributes}
%=============================================================

%-------------------------------------------------------------
\subsection{Reliability}
%-------------------------------------------------------------

The RAG engine and ingestion pipeline are initialised as module-level singletons in
\texttt{rag\_service.py}, loaded once when the FastAPI server starts. If the Groq API
is temporarily unavailable, the backend logs the error and returns a structured HTTP 500
rather than taking down the server process. MongoDB is handled the same way. The guiding
principle was straightforward: one broken external dependency should not kill the whole
application, and a structured error message is almost always more useful than a crash.

%-------------------------------------------------------------
\subsection{Availability}
%-------------------------------------------------------------

The system is packaged with Docker Compose. A single \texttt{docker-compose up} command
starts the FastAPI backend, the React frontend and MongoDB together. Once running, the
application is available at \texttt{localhost:5173} (frontend) and \texttt{localhost:8000}
(API). No further configuration is required beyond setting the environment variables
described in the CLAUDE.md file.

%-------------------------------------------------------------
\subsection{Security}
%-------------------------------------------------------------

Passwords go through bcrypt before they ever touch the database --- there is no plaintext
stored anywhere in the system. Access tokens live exclusively in Zustand in-memory state
and disappear when the browser tab closes. Refresh tokens go into HTTP-only cookies that
JavaScript cannot read, which closes off the most common XSS token-theft vector. CORS is
configured in \texttt{main.py} to restrict cross-origin requests to the deployed frontend
origin; wildcard origins are not used.

%-------------------------------------------------------------
\subsection{Maintainability}
%-------------------------------------------------------------

The backend splits cleanly into independent FastAPI routers, one file each for auth,
sessions and query/pipeline. Adding a new API surface means writing one router file and
one \texttt{app.include\_router()} call --- nothing else changes. On the frontend, all
data-fetching and auth logic lives in custom hooks (\texttt{useAuth}, \texttt{useChat},
\texttt{useSessions}), keeping UI components stateless. That separation makes it
much easier to test business logic independently of the rendering layer.

%-------------------------------------------------------------
\subsection{Portability}
%-------------------------------------------------------------

Everything runs inside Docker containers. The host machine needs Docker Engine and nothing
else. Python 3.11 and Node 20 are fixed in the build image specifications so that a build
on one machine produces the same result as a build on any other --- the
``works on my machine'' problem is eliminated by design.

%-------------------------------------------------------------
\subsection{Performance}
%-------------------------------------------------------------

Heavy ML imports --- \texttt{sentence-transformers} and \texttt{chromadb} --- are deferred
inside \texttt{init\_rag()} and loaded once at startup rather than at module import time.
This keeps FastAPI's startup fast and avoids reloading the embedding model on every
request. The \texttt{all-MiniLM-L6-v2} model produces 384-dimensional embeddings; at that
dimensionality, a ChromaDB cosine search across several thousand chunks runs in well under
100\,ms on a standard CPU. That leaves the Groq API call as the primary latency variable,
not the retrieval step.

%=============================================================
\section{Database Requirements}
%=============================================================

MediLingo AI uses two persistence layers.

\textbf{MongoDB} stores three collections:
\begin{itemize}
  \item \texttt{users} --- one document per registered account, containing the UUID4
        string ID, display name, email, bcrypt password hash, role and active flag
  \item \texttt{sessions} --- one document per chat session, containing the session ID,
        the owning user's ID and a display title
  \item \texttt{messages} --- one document per message (user query or assistant response),
        containing the parent session ID, the message role
        (\texttt{user} or \texttt{assistant}), the content text, the language code and a
        timestamp
\end{itemize}

\textbf{ChromaDB} stores one collection (\texttt{medical\_documents}):
\begin{itemize}
  \item Each record holds the chunk text, its 384-dimensional embedding vector, and
        metadata including the source filename and chunk index
  \item The collection is populated by the ingestion pipeline and queried at inference
        time using cosine similarity
  \item The collection persists to \texttt{backend/vectorstore/} on the host filesystem
\end{itemize}

%=============================================================
\section{Use Cases}
%=============================================================

%-------------------------------------------------------------
\subsection{Use Cases List}
%-------------------------------------------------------------

\renewcommand{\arraystretch}{1.3}
\begin{table}[H]
\centering
\small
\caption{User Use Cases}
\label{tab:user_uc}
\begin{tabular}{|>{\bfseries\color{navyblue}\centering}p{1.5cm}|p{4.5cm}|p{7cm}|}
\hline
\theader{UC ID} & \theader{Use Case Name} & \theader{Brief Description} \\
\hline
UC-1 & Register User Account    & Create a new MediLingo AI account with name, email and password. \\ \hline
UC-2 & Login                    & Authenticate an existing account and receive JWT tokens. \\ \hline
UC-3 & Submit Medical Query     & Send a multilingual medical question and receive a RAG-grounded response. \\ \hline
UC-4 & Manage Chat Sessions     & Create, continue, view and delete persistent chat sessions. \\ \hline
UC-5 & Use Voice Input          & Dictate a query via microphone using the Web Speech API. \\ \hline
UC-6 & Use Voice Output         & Have an assistant response read aloud via text-to-speech. \\ \hline
UC-7 & Manage Profile           & View and update display name and email address. \\ \hline
UC-8 & Reset Password           & Recover account access via an email one-time passcode. \\ \hline
\end{tabular}
\end{table}

\begin{table}[H]
\centering
\small
\caption{Admin Use Cases}
\label{tab:admin_uc}
\begin{tabular}{|>{\bfseries\color{navyblue}\centering}p{1.5cm}|p{4.5cm}|p{7cm}|}
\hline
\theader{UC ID} & \theader{Use Case Name} & \theader{Brief Description} \\
\hline
UC-9  & Manage Users            & View all accounts; activate or deactivate individual users. \\ \hline
UC-10 & Ingest Medical Documents & Trigger PDF extraction, chunking, embedding and ChromaDB indexing. \\ \hline
\end{tabular}
\end{table}

%-------------------------------------------------------------
\subsection{Use Case Diagram}
%-------------------------------------------------------------

\begin{figure}[H]
  \centering
  \includegraphics[width=0.85 \textwidth]{use_case_diagram}
  \caption{MediLingo AI --- Use Case Diagram}
  \label{fig:ucd}
\end{figure}

%-------------------------------------------------------------
\subsection{Use Cases Description}
%-------------------------------------------------------------

\renewcommand{\arraystretch}{1.4}

% ─── UC-1 ───────────────────────────────────────────────────
\begin{table}[H]
\centering\small
\caption{UC-1: Register User Account}
\label{tab:uc1}
\begin{tabular}{|>{\bfseries\color{navyblue}}p{3.2cm}|p{9.8cm}|}
\hline
\multicolumn{2}{|c|}{\cellcolor{navyblue}\textbf{\color{white}UC-1: Register User Account}} \\ \hline
Use Case ID     & UC-1 \\ \hline
Actor           & Unregistered User \\ \hline
Description     & A new user creates a MediLingo AI account by supplying a display name, email address and password. \\ \hline
Preconditions   & The provided email address is not already associated with an existing account. \\ \hline
Main Flow       &
  1.\ User navigates to the Register page. \newline
  2.\ User enters display name, email address and password. \newline
  3.\ System validates that the email is not already in use. \newline
  4.\ System hashes the password with bcrypt and creates a UserDoc in MongoDB. \newline
  5.\ System issues a 15-minute access token and sets a 7-day HTTP-only refresh cookie. \newline
  6.\ User is redirected to the chat interface. \\ \hline
Alternate Flow  &
  \textbf{A1} -- Email already registered: system returns HTTP~400 and the user corrects the email field. \newline
  \textbf{A2} -- Password below minimum length: the form is rejected client-side before submission. \\ \hline
Postconditions  & A UserDoc exists in the \texttt{users} collection. The user is authenticated with valid JWT tokens. \\ \hline
\end{tabular}
\end{table}

% ─── UC-2 ───────────────────────────────────────────────────
\begin{table}[H]
\centering\small
\caption{UC-2: Login}
\label{tab:uc2}
\begin{tabular}{|>{\bfseries\color{navyblue}}p{3.2cm}|p{9.8cm}|}
\hline
\multicolumn{2}{|c|}{\cellcolor{navyblue}\textbf{\color{white}UC-2: Login}} \\ \hline
Use Case ID     & UC-2 \\ \hline
Actor           & Registered User / Admin User \\ \hline
Description     & An existing user authenticates with their email and password to receive JWT tokens. \\ \hline
Preconditions   & A valid account exists for the provided email. The account is active. \\ \hline
Main Flow       &
  1.\ User navigates to the Login page. \newline
  2.\ User enters email address and password. \newline
  3.\ System verifies the bcrypt hash against the stored value. \newline
  4.\ System issues a new access token and sets a new refresh cookie. \newline
  5.\ User is redirected to the chat interface. \\ \hline
Alternate Flow  &
  \textbf{A1} -- Incorrect credentials: system returns HTTP~401 and the user is prompted to retry. \newline
  \textbf{A2} -- Account deactivated: system returns HTTP~403 with an appropriate error message. \\ \hline
Postconditions  & The user's access token is held in Zustand memory. A refresh cookie is set in the browser. \\ \hline
\end{tabular}
\end{table}

% ─── UC-3 ───────────────────────────────────────────────────
\begin{table}[H]
\centering\small
\caption{UC-3: Submit Medical Query}
\label{tab:uc3}
\begin{tabular}{|>{\bfseries\color{navyblue}}p{3.2cm}|p{9.8cm}|}
\hline
\multicolumn{2}{|c|}{\cellcolor{navyblue}\textbf{\color{white}UC-3: Submit Medical Query}} \\ \hline
Use Case ID     & UC-3 \\ \hline
Actor           & Authenticated User \\ \hline
Description     & The user submits a medical question in their chosen language and receives a RAG-grounded response in the same language. \\ \hline
Preconditions   & User is authenticated. The RAG engine is initialised. At least one document has been ingested into ChromaDB. \\ \hline
Main Flow       &
  1.\ User selects a language (English / Urdu / Roman Urdu) and types a query. \newline
  2.\ Frontend prepends the language instruction to the query text. \newline
  3.\ Frontend sends POST \texttt{/query} with the prefixed query and session ID. \newline
  4.\ Backend strips the language prefix; the clean query is embedded with \texttt{all-MiniLM-L6-v2}. \newline
  5.\ Backend performs cosine similarity search in ChromaDB and retrieves the top-K passages. \newline
  6.\ Backend builds a prompt containing the language instruction, retrieved passages and query. \newline
  7.\ Groq LLaMA~3.3-70b generates a response in the selected language. \newline
  8.\ Backend saves user message and assistant response as MessageDocs in MongoDB. \newline
  9.\ Response is returned to the frontend and displayed in the chat interface. \\ \hline
Alternate Flow  &
  \textbf{A1} -- Groq API unavailable: system returns HTTP~500 with ``Service temporarily unavailable''. \newline
  \textbf{A2} -- No relevant passages found: LLM is prompted to indicate insufficient information in the knowledge base. \newline
  \textbf{A3} -- Access token expired: \texttt{authenticatedFetch} silently calls \texttt{/auth/refresh} and retries. \\ \hline
Postconditions  & The query and response are saved to the session. The chat view is updated. \\ \hline
\end{tabular}
\end{table}

% ─── UC-4 ───────────────────────────────────────────────────
\begin{table}[H]
\centering\small
\caption{UC-4: Manage Chat Sessions}
\label{tab:uc4}
\begin{tabular}{|>{\bfseries\color{navyblue}}p{3.2cm}|p{9.8cm}|}
\hline
\multicolumn{2}{|c|}{\cellcolor{navyblue}\textbf{\color{white}UC-4: Manage Chat Sessions}} \\ \hline
Use Case ID     & UC-4 \\ \hline
Actor           & Authenticated User \\ \hline
Description     & The user creates, switches between and deletes named chat sessions via the session sidebar. \\ \hline
Preconditions   & User is authenticated. \\ \hline
Main Flow       &
  1.\ User clicks ``New Session'' in the sidebar; a new SessionDoc is created in MongoDB. \newline
  2.\ User clicks an existing session to load its message history. \newline
  3.\ User deletes a session; the SessionDoc and all child MessageDocs are removed. \\ \hline
Alternate Flow  &
  \textbf{A1} -- Network error during session creation: the sidebar shows an error and the session is not added. \\ \hline
Postconditions  & The selected session is active and its messages are loaded in the chat area. \\ \hline
\end{tabular}
\end{table}

% ─── UC-5 ───────────────────────────────────────────────────
\begin{table}[H]
\centering\small
\caption{UC-5: Use Voice Input}
\label{tab:uc5}
\begin{tabular}{|>{\bfseries\color{navyblue}}p{3.2cm}|p{9.8cm}|}
\hline
\multicolumn{2}{|c|}{\cellcolor{navyblue}\textbf{\color{white}UC-5: Use Voice Input}} \\ \hline
Use Case ID     & UC-5 \\ \hline
Actor           & Authenticated User \\ \hline
Description     & The user dictates their query via the microphone button; the transcript populates the query input field. \\ \hline
Preconditions   & User is authenticated. Browser supports the Web Speech API. Microphone permission has been granted. \\ \hline
Main Flow       &
  1.\ User clicks the microphone button in the chat interface. \newline
  2.\ Browser requests microphone permission if not already granted. \newline
  3.\ Speech recognition starts; the user speaks their query. \newline
  4.\ Transcript text is placed in the query input field. \newline
  5.\ User submits the query normally (UC-3). \\ \hline
Alternate Flow  &
  \textbf{A1} -- Microphone permission denied: an error alert is shown and recognition does not start. \newline
  \textbf{A2} -- Browser does not support Web Speech API: the microphone button is hidden. \\ \hline
Postconditions  & Query input field contains the transcribed text, ready for submission. \\ \hline
\end{tabular}
\end{table}

% ─── UC-6 ───────────────────────────────────────────────────
\begin{table}[H]
\centering\small
\caption{UC-6: Use Voice Output}
\label{tab:uc6}
\begin{tabular}{|>{\bfseries\color{navyblue}}p{3.2cm}|p{9.8cm}|}
\hline
\multicolumn{2}{|c|}{\cellcolor{navyblue}\textbf{\color{white}UC-6: Use Voice Output}} \\ \hline
Use Case ID     & UC-6 \\ \hline
Actor           & Authenticated User \\ \hline
Description     & After an assistant response is displayed, the system reads it aloud in the user's selected language using the browser's SpeechSynthesis API. \\ \hline
Preconditions   & A response has been received and rendered. Browser supports SpeechSynthesis. \\ \hline
Main Flow       &
  1.\ Assistant response is received by the frontend. \newline
  2.\ \texttt{useVoiceOutput} hook passes the response text to the SpeechSynthesis API. \newline
  3.\ Browser reads the text aloud in the configured language voice. \\ \hline
Alternate Flow  &
  \textbf{A1} -- SpeechSynthesis not available: response is displayed as text only; no error is shown. \\ \hline
Postconditions  & The response text has been read aloud or silently displayed if TTS is unavailable. \\ \hline
\end{tabular}
\end{table}

% ─── UC-7 ───────────────────────────────────────────────────
\begin{table}[H]
\centering\small
\caption{UC-7: Manage Profile}
\label{tab:uc7}
\begin{tabular}{|>{\bfseries\color{navyblue}}p{3.2cm}|p{9.8cm}|}
\hline
\multicolumn{2}{|c|}{\cellcolor{navyblue}\textbf{\color{white}UC-7: Manage Profile}} \\ \hline
Use Case ID     & UC-7 \\ \hline
Actor           & Authenticated User \\ \hline
Description     & The user views and updates their display name or email address from the profile page. \\ \hline
Preconditions   & User is authenticated. \\ \hline
Main Flow       &
  1.\ User navigates to the Profile page. \newline
  2.\ Current display name and email are pre-filled in the form. \newline
  3.\ User edits one or both fields and clicks Save. \newline
  4.\ System updates the UserDoc in MongoDB and returns the updated values. \\ \hline
Alternate Flow  &
  \textbf{A1} -- New email already taken: system returns HTTP~400 and the user corrects input. \\ \hline
Postconditions  & The UserDoc in MongoDB reflects the updated display name or email. \\ \hline
\end{tabular}
\end{table}

% ─── UC-8 ───────────────────────────────────────────────────
\begin{table}[H]
\centering\small
\caption{UC-8: Reset Password}
\label{tab:uc8}
\begin{tabular}{|>{\bfseries\color{navyblue}}p{3.2cm}|p{9.8cm}|}
\hline
\multicolumn{2}{|c|}{\cellcolor{navyblue}\textbf{\color{white}UC-8: Reset Password}} \\ \hline
Use Case ID     & UC-8 \\ \hline
Actor           & Registered User \\ \hline
Description     & A user who has forgotten their password recovers account access via a one-time passcode sent to their registered email address. \\ \hline
Preconditions   & An account exists for the email address provided. \\ \hline
Main Flow       &
  1.\ User navigates to Forgot Password and enters their registered email. \newline
  2.\ System generates a time-limited OTP and sends it via the email service. \newline
  3.\ User enters the OTP on the Reset Password page. \newline
  4.\ User sets and confirms a new password. \newline
  5.\ System hashes the new password and updates the UserDoc. \newline
  6.\ User is redirected to the Login page. \\ \hline
Alternate Flow  &
  \textbf{A1} -- Email not registered: system returns HTTP~404; no OTP is sent. \newline
  \textbf{A2} -- OTP expired: user must restart the process from step 1. \newline
  \textbf{A3} -- Passwords do not match: form is rejected client-side before submission. \\ \hline
Postconditions  & The user's password hash in MongoDB is updated. The user can log in with the new password. \\ \hline
\end{tabular}
\end{table}

% ─── UC-9 ───────────────────────────────────────────────────
\begin{table}[H]
\centering\small
\caption{UC-9: Manage Users (Admin)}
\label{tab:uc9}
\begin{tabular}{|>{\bfseries\color{navyblue}}p{3.2cm}|p{9.8cm}|}
\hline
\multicolumn{2}{|c|}{\cellcolor{navyblue}\textbf{\color{white}UC-9: Manage Users (Admin)}} \\ \hline
Use Case ID     & UC-9 \\ \hline
Actor           & Admin User \\ \hline
Description     & An admin views the full user list and can activate or deactivate individual accounts. \\ \hline
Preconditions   & Admin is authenticated via the admin login route. \\ \hline
Main Flow       &
  1.\ Admin navigates to the Admin Panel. \newline
  2.\ System displays a paginated list of all registered users with their status. \newline
  3.\ Admin clicks the toggle next to a user to activate or deactivate their account. \newline
  4.\ System updates the \texttt{is\_active} flag in the UserDoc. \\ \hline
Alternate Flow  &
  \textbf{A1} -- Admin attempts to deactivate their own account: system returns an error to prevent lockout. \\ \hline
Postconditions  & The targeted user's active status is updated. If deactivated, their next login attempt returns HTTP~403. \\ \hline
\end{tabular}
\end{table}

% ─── UC-10 ───────────────────────────────────────────────────
\begin{table}[H]
\centering\small
\caption{UC-10: Ingest Medical Documents (Admin)}
\label{tab:uc10}
\begin{tabular}{|>{\bfseries\color{navyblue}}p{3.2cm}|p{9.8cm}|}
\hline
\multicolumn{2}{|c|}{\cellcolor{navyblue}\textbf{\color{white}UC-10: Ingest Medical Documents (Admin)}} \\ \hline
Use Case ID     & UC-10 \\ \hline
Actor           & Admin User \\ \hline
Description     & An admin triggers the RAG ingestion pipeline to extract, chunk, embed and index PDF files from the datasets directory into ChromaDB. \\ \hline
Preconditions   & Admin is authenticated. One or more PDF files are present in \texttt{backend/datasets/}. \\ \hline
Main Flow       &
  1.\ Admin navigates to the Admin Panel and clicks ``Ingest Documents''. \newline
  2.\ Frontend sends POST \texttt{/pipeline/ingest-sync}. \newline
  3.\ Backend reads all PDF files from \texttt{backend/datasets/} using \texttt{pdfplumber}. \newline
  4.\ Extracted text is split into 1000-character chunks with 200-character overlap. \newline
  5.\ Each chunk is embedded with \texttt{all-MiniLM-L6-v2}. \newline
  6.\ Embeddings and metadata are stored in the \texttt{medical\_documents} ChromaDB collection. \newline
  7.\ System returns a summary of ingested documents and chunk count. \\ \hline
Alternate Flow  &
  \textbf{A1} -- No PDF files found: system returns HTTP~400 with ``No documents found in datasets directory''. \newline
  \textbf{A2} -- PDF is password-protected or corrupted: the file is skipped and a warning is included in the response. \\ \hline
Postconditions  & ChromaDB contains embeddings for all successfully processed chunks. Subsequent queries can retrieve passages from the newly ingested documents. \\ \hline
\end{tabular}
\end{table}

%=============================================================
\section{System Sequence Diagrams}
%=============================================================

A System Sequence Diagram (SSD) shows the interaction between an actor and the system
boundary for a given use case. The system is a black box at this level of abstraction ---
exactly what the SRS perspective calls for. The internals --- how FastAPI coordinates
the RAG engine, ChromaDB and MongoDB to produce a response --- are covered in the detailed
sequence diagrams of Chapter~3.

%-------------------------------------------------------------
\subsection{SSD-1: Register User Account}
%-------------------------------------------------------------
\begin{figure}[H]
  \centering
  \includegraphics[width=0.85\textwidth]{ssd1_register}
  \caption{SSD-1: Register User Account}
  \label{fig:ssd1}
\end{figure}

%-------------------------------------------------------------
\subsection{SSD-2: Login}
%-------------------------------------------------------------
\begin{figure}[H]
  \centering
  \includegraphics[width=0.85\textwidth]{ssd2_login}
  \caption{SSD-2: Login}
  \label{fig:ssd2}
\end{figure}

%-------------------------------------------------------------
\subsection{SSD-3: Submit Medical Query}
%-------------------------------------------------------------
\begin{figure}[H]
  \centering
  \includegraphics[width=\textwidth]{ssd3_submit_query}
  \caption{SSD-3: Submit Medical Query}
  \label{fig:ssd3}
\end{figure}

%-------------------------------------------------------------
\subsection{SSD-4: Manage Chat Sessions}
%-------------------------------------------------------------
\begin{figure}[H]
  \centering
  \includegraphics[width=0.90\textwidth]{ssd4_sessions}
  \caption{SSD-4: Manage Chat Sessions (Create / Load / Delete)}
  \label{fig:ssd4}
\end{figure}

%-------------------------------------------------------------
\subsection{SSD-5: Use Voice Input}
%-------------------------------------------------------------
\begin{figure}[H]
  \centering
  \includegraphics[width=0.85\textwidth]{ssd5_voice_input}
  \caption{SSD-5: Use Voice Input}
  \label{fig:ssd5}
\end{figure}

%-------------------------------------------------------------
\subsection{SSD-6: Use Voice Output}
%-------------------------------------------------------------
\begin{figure}[H]
  \centering
  \includegraphics[width=0.75\textwidth]{ssd6_voice_output}
  \caption{SSD-6: Use Voice Output}
  \label{fig:ssd6}
\end{figure}

%-------------------------------------------------------------
\subsection{SSD-7: Manage Profile}
%-------------------------------------------------------------
\begin{figure}[H]
  \centering
  \includegraphics[width=0.85\textwidth]{ssd7_profile}
  \caption{SSD-7: Manage Profile}
  \label{fig:ssd7}
\end{figure}

%-------------------------------------------------------------
\subsection{SSD-8: Reset Password}
%-------------------------------------------------------------
\begin{figure}[H]
  \centering
  \includegraphics[width=0.90\textwidth]{ssd8_reset_password}
  \caption{SSD-8: Reset Password via Email OTP}
  \label{fig:ssd8}
\end{figure}

%-------------------------------------------------------------
\subsection{SSD-9: Manage Users (Admin)}
%-------------------------------------------------------------
\begin{figure}[H]
  \centering
  \includegraphics[width=0.85\textwidth]{ssd9_manage_users}
  \caption{SSD-9: Manage Users (Admin)}
  \label{fig:ssd9}
\end{figure}

%-------------------------------------------------------------
\subsection{SSD-10: Ingest Medical Documents (Admin)}
%-------------------------------------------------------------
\begin{figure}[H]
  \centering
  \includegraphics[width=0.90\textwidth]{ssd10_ingest}
  \caption{SSD-10: Ingest Medical Documents (Admin)}
  \label{fig:ssd10}
\end{figure}

%=============================================================
\section{Domain Model}
%=============================================================

The domain model in Figure~\ref{fig:domain} shows the three core persistent entities and
how they relate. The structure is deliberately simple. A \textbf{User} owns zero or more
\textbf{Sessions}. Each \textbf{Session} contains zero or more \textbf{Messages}. When a
session is deleted, everything under it goes with it --- there are no orphaned messages
left behind in the database. All three entities are stored in MongoDB; messages are
embedded inside their parent session document rather than in a separate collection.

\begin{figure}[H]
  \centering
  \includegraphics[width=\textwidth]{domain_model}
  \caption{MediLingo AI --- Domain Model}
  \label{fig:domain}
\end{figure}
